% ============================================================
%  Chapter 2: Related Work
%  File: chapters/02_related_work.tex
%  Bibliography: chapters/refs_02_related.bib
% ============================================================

\chapter{Related Work}
\label{ch:related}

Our work sits at the intersection of six research traditions:
neuroevolution, growing and self-organizing networks, modular neural
networks, sparse connectivity, digital circuit synthesis with neural
networks, and universal computation. This chapter reviews each in turn
and identifies the gap that \modnet{} addresses.

% ------------------------------------------------------------
\section{Neuroevolution: Evolving Network Topologies}
\label{sec:related-neuroevolution}
% ------------------------------------------------------------

The idea that a neural network's topology can be evolved rather than
designed dates back to the early 1990s. \citet{gruau1994neural}
introduced \emph{Cellular Encoding}, in which a genetic program
simultaneously evolves the architecture, weights, thresholds, and biases
of a neural network. The genotype is a tree of graph-rewriting rules
inspired by L-systems; through development, this tree unfolds into a
network of arbitrary size and shape. Gruau demonstrated that cellular
encoding could generate modular neural networks and used it to evolve a
controller for a hexapod robot. This early work established the principle
that complex structure can emerge from a compact generative description.

The NeuroEvolution of Augmenting Topologies (NEAT) algorithm of
\citet{stanley2002evolving} became the canonical framework for
topology-evolving neuroevolution. NEAT starts from a minimal network and
incrementally adds nodes and connections through structural mutation. It
introduces three key ideas:

\begin{enumerate}
  \item \textbf{Historical markings} to allow meaningful crossover
    between networks of different topologies.
  \item \textbf{Speciation} to protect structural innovation from
    premature competition.
  \item \textbf{Complexification} from small initial topologies.
\end{enumerate}

NEAT was shown to outperform fixed-topology methods on reinforcement
learning benchmarks \citep{stanley2002efficient} and has since been
applied to domains ranging from autonomous rover navigation
\citep{shrestha2025reinforced} to scheduling problems
\citep{neat2024scheduling}.

HyperNEAT \citep{stanley2009hypercube} extended NEAT with an indirect
encoding based on Compositional Pattern-Producing Networks (CPPNs).
Instead of directly encoding each connection, HyperNEAT evolves a CPPN
whose output at any pair of coordinates specifies the weight of the
corresponding connection. This geometric encoding exploits regularity in
the substrate and scales to networks with millions of connections
\citep{clune2011performance}. The Enhanced Hypercube-based Encoding
(ES-HyperNEAT) of \citet{risi2012enhanced} further allowed the placement,
density, and connectivity of neurons to be evolved rather than fixed.

Recent work has focused on accelerating NEAT and extending its reach.
\citet{wang2024tensorized} introduced a tensorization method that
transforms the diverse topologies of NEAT into uniformly shaped tensors,
enabling GPU acceleration and reducing training time by an order of
magnitude. \citet{lyubchev2025revisiting} revisited the original NEAT
formalism and proposed improvements to the mutation and selection
operators.

Despite their success, these methods share a common assumption: the
network is organized into layers or a predefined substrate geometry.
HyperNEAT's CPPN operates on a grid; NEAT's complexification adds nodes
to hidden layers; tensorized NEAT pads topologies into fixed-shape
tensors. None of these approaches explores fully layer-free computation
in which \emph{any node may connect to any other node}, including
itself, and in which there is no modular scaffold.

% ------------------------------------------------------------
\section{Growing and Self-Organizing Networks}
\label{sec:related-growing}
% ------------------------------------------------------------

An alternative to evolving a fixed topology is to grow the network
through local rules, either developmental or activity-dependent.

\subsection{Constructive Growth}

The Cascade-Correlation architecture of \citet{fahlman1990cascade}
begins with a minimal network and adds hidden units one by one, training
each new unit to maximize the correlation between its output and the
residual error. The Growing Neural Gas algorithm of
\citet{fritzke1995growing} uses competitive Hebbian learning to build
topological representations of input space incrementally. Both
approaches demonstrate the power of constructive growth, but neither
addresses the problem of recurrence, nor the possibility of a fully
layer-free topology.

\subsection{Developmental Systems}

A related line of research investigates neural networks inspired by
biological development. Rather than evolving a fixed structure, these
systems grow, prune, and reorganize themselves through local rules.

\citet{plantec2024evolving} proposed the Lifelong Neural Developmental
Program (LNDP), a class of self-organizing networks capable of synaptic
and structural plasticity in an activity- and reward-dependent manner.
Their results demonstrated that structural plasticity is advantageous in
environments requiring fast adaptation or with non-stationary rewards.
The LNDP grows from a single cell using a Neural Developmental Program,
but the growth rules themselves are evolved, not the topology directly.

The GRN-SO-RC model of \citet{hajizada2024developmental} uses a gene
regulatory network (GRN) to regulate synaptic and structural plasticity
in a reservoir computing framework. The structure of the reservoir
self-organizes to adapt to the specific task through the GRN mechanism.
This work is notable for demonstrating that self-organized reservoir
structures can match or exceed the performance of hand-designed
topologies.

\citet{hazan2024self} proposed a biologically inspired developmental
algorithm that grows a functional, layered neural network from a single
initial cell. The growth process is driven by noise and local
self-organization, and the resulting networks achieve competitive
performance on standard benchmarks.

The broader relationship between neuroevolution and gene regulatory
networks was reviewed by \citet{gerges2025neuroevolution}, who noted
that GRNs are more robust to mutation and thus more evolvable than direct
neuroevolution. This observation motivates the use of compact, generative
encodings rather than direct encodings of network structure.

\subsection{Neural Cellular Automata}

Neural cellular automata, introduced by \citet{mordvintsev2020growing},
train a cellular automaton to grow and regenerate a target image from a
single seed. \citet{bena2025path} recently demonstrated that neural
cellular automata can achieve universal computation in a continuous
setting.

\modnet{} shares with neural cellular automata the idea that computation
can be local, recurrent, and free of any predefined layer structure.
However, while neural cellular automata operate on a fixed grid with a
shared update rule, \modnet{} allows the topology itself---including the
number and connectivity of nodes---to be evolved.

% ------------------------------------------------------------
\section{Modular Neural Networks}
\label{sec:related-modular}
% ------------------------------------------------------------

Modularity is a well-established design principle in software
engineering. \citet{parnas1972modularity} argued that systems should be
decomposed into modules with clean interfaces, hidden implementation
details, low coupling between modules, and high cohesion within modules.
This principle has been remarkably successful in software design
\citep{larman2004applying}, and it is natural to ask whether it applies
to neural computation as well.

\citet{alho2024modularity} conducted a systematic review of $86$ studies
on modular neural networks and found that nearly two-thirds report
improvements in task accuracy compared to monolithic architectures, with
$82\%$ reporting benefits in training or inference efficiency. This
suggests that modularity is not merely a design preference but a
computationally meaningful structural property.

In the context of neuroevolution, \citet{khare2005coevolutionary}
proposed a co-evolutionary model in which a population of modules is
evolved alongside a population of complete systems that combine modules.
They demonstrated that if a particular task decomposition is beneficial,
it can be discovered through this mechanism. However, their model
requires the designer to specify the number of modules and the interface
between them.

\citet{landassuri2012evolution} investigated the emergence of modularity
in evolved neural networks and found that pure-modular architectures
emerge at low connectivity values, and that the more different two tasks
are, the greater the modularity obtained during evolution. This suggests
that modularity is an emergent property of evolutionary search under
certain conditions, not a design choice imposed from above.

More recent work has examined modularity in the context of reinforcement
learning. \citet{munn2024towards} used NEAT to generate networks of
varying modularity on three RL tasks and found that structural
modularity correlated with performance on tasks that admitted natural
task decomposition, but not on others.

Related work in hierarchical modular networks includes
\citet{rodriguez2025hierarchical}, who showed that multi-level
hierarchical modular networks, when implemented as recurrent neural
network reservoirs, enhance memory capacity, support multitasking, and
give rise to a broader range of temporal dynamics. The modular growth
curriculum of \citet{modulargrowth2024} similarly showed that gradual
modular growth of RNNs provides advantages for learning increasingly
complex tasks.

In earlier work, we studied a \emph{modular} variant of \modnet{},
where nodes were grouped into a fixed number of modules. The present
work is fully \emph{layer-free}: there are no modules at all. This is a
deliberate departure. Our results show that the layer-free variant is
equal or superior on every benchmark, including a decisive win on the
multiplication task, where it uses only $13$ nodes versus $18$ in the
modular variant. We discuss this comparison in detail in
Chapter~\ref{ch:discussion}.

% ------------------------------------------------------------
\section{Sparse Connectivity}
\label{sec:related-sparse}
% ------------------------------------------------------------

Biological neural networks are highly sparse: the probability that any
two neurons are connected is on the order of $10^{-6}$
\citep{lecun2015deep}. This sparsity is not accidental; it reflects
metabolic and wiring-length constraints \citep{cherniak1994component,
chklovskii2004wiring}.

In artificial neural networks, sparsity has been studied primarily as a
compression technique. The Sparse Evolutionary Training (SET) method of
\citet{mocanu2018scalable} evolves an initial sparse topology into a
scale-free topology during training, allowing the training of networks
with over one million neurons on commodity hardware.
\citet{liu2019sparse} extended this approach to show that vanilla
SET-MLP can be implemented from scratch using just pure Python and
SciPy, and that the resulting networks achieve competitive accuracy on
standard benchmarks. \citet{cosinesimilarity2019} proposed a
cosine-similarity-based selection criterion for SET that improves the
efficiency of the evolutionary topology update.

More recent work has considered pruning from an evolutionary
perspective. \citet{shah2026pruning} modeled parameter groups as
populations whose influence evolves under selection pressure, arguing
that sparsity is not merely a design choice but an emergent property of
selection dynamics under resource constraints.
\citet{hershey2025balancing} showed that sparse RNNs can be balanced
with appropriate hyperparameterization, and that this balance benefits
meta-learning.

The relationship between sparsity and modularity is intimate.
\citet{landassuri2012evolution} found that pure-modular architectures
emerge at low connectivity values, suggesting that sparsity and
modularity are two faces of the same underlying structural pressure. In
\modnet{}, the pruning mechanism operates without any explicit sparsity
target: nodes and connections are removed when they become inactive or
isolated, and the resulting networks end up with $27\%$ to $67\%$
active-wire density depending on the task.

% ------------------------------------------------------------
\section{Neural Networks for Digital Circuits}
\label{sec:related-circuits}
% ------------------------------------------------------------

The idea that neural networks can implement digital logic has a long
history. The perceptron \citep{rosenblatt1958perceptron} was originally
conceived as a threshold unit, and the connection between threshold
circuits and neural networks was formalized early
\citep{minsky1969perceptrons}.

\citet{hajek1990adder} and others demonstrated that simple neural
networks can learn to compute arithmetic operations such as addition and
multiplication. \citet{siu1991adder} showed that a three-layer
feedforward network with threshold units can compute binary addition
exactly. The connection between neural networks and Boolean circuits was
further explored by \citet{siu1995discrete}, who argued that
threshold circuits are the natural computational model for neural
networks.

For sorting networks, the problem of learning a sorting algorithm with
a neural network was studied by \citet{knuth1998sorting}, who noted
that sorting is a canonical example of a problem that requires
\emph{non-monotone} Boolean logic. A sorting network can be implemented
as a circuit of comparators, and learning such a circuit from scratch
is nontrivial.

In the modern era, the focus has shifted to \emph{learning} digital
circuits with differentiable relaxations. \citet{petersen2022deep}
proposed differentiable sorting networks based on the
Gumbel-Sinkhorn operator. \citet{grover2019differentiable} introduced
NeuralSort, a continuous relaxation of sorting. These approaches are
elegant, but they rely on gradient information and a predefined
architectural template.

\modnet{} approaches the problem differently: rather than relaxing
sorting into a differentiable function, we evolve a Boolean circuit
from scratch. To the best of our knowledge, this is the first
demonstration of a fully layer-free evolved network that solves full
adder, 2-bit adder, comparator, 4-to-1 multiplexer, and 4-bit sorting
network, all with exact $100\%$ accuracy.

% ------------------------------------------------------------
\section{Universal Computation}
\label{sec:related-universal}
% ------------------------------------------------------------

The theoretical foundations of neural computation were established by
\citet{siegelmann1995computational}, who proved that recurrent neural
networks with rational weights are Turing complete. More recently,
\citet{perez2021turing} showed that transformers are Turing complete
without external memory, and \citet{bates2026pytorch} released a
library of Turing-complete neural networks for reproducible
experimentation. \citet{stogin2024provably} constructed a provably
stable neural network Turing machine with finite precision and time.

Cellular automata provide a complementary perspective. Rule 110 was
proven Turing complete by \citet{cook2004universality}, and reversible
two-dimensional cellular spaces were shown to support universal
computation by \citet{morita2002universal}.

\modnet{} is, in principle, Turing complete: because it permits
arbitrary recurrent connections and can grow to arbitrary size, it
contains recurrent neural networks as a special case. We do not prove
this formally here, but we note it as a natural extension of our work.

% ------------------------------------------------------------
\section{Biological Inspiration: Criticality and Small-World Topology}
\label{sec:related-biology}
% ------------------------------------------------------------

Biological neural networks exhibit two structural properties that have
been extensively studied: critical dynamics and small-world topology.

The criticality hypothesis states that the brain operates near the phase
transition between order and disorder, where information transmission is
maximized \citep{beggs2003neuronal}. \citet{poil2025network} reviewed
the mechanisms by which this critical state is maintained, identifying
self-organized criticality (SOC) as a leading candidate.
\citet{organoid2026criticality} demonstrated that human forebrain
organoids self-organize to a computationally favorable critical state
without external input. In the context of modular networks,
\citet{hierarchicalmodular2025} showed that structural modularity tunes
mesoscale criticality in biological neuronal networks.

Small-world topology---characterized by high clustering and short path
lengths---has also been studied extensively. \citet{emergence2024} and
\citet{evolutionarily2026} showed that spontaneously evolved or
evolutionarily optimized topologies exhibit biologically plausible
small-world properties, including densely connected hub nodes, short
paths, and long-tailed degree distributions. \citet{topology2004}
showed that scale-free topologies store and retrieve binary patterns
more efficiently than random-diluted networks with the same number of
synapses.

These biological and topological findings provide both motivation and
validation for \modnet{}. Our architecture does not enforce small-world
or scale-free structure; we observe that sparsity, cross-scale
integration, and recurrence emerge from evolutionary pressure alone,
without explicit regularization for any of these properties.

% ------------------------------------------------------------
\section{Neural Architecture Search}
\label{sec:related-nas}
% ------------------------------------------------------------

Neural Architecture Search (NAS) automates the design of network
architectures. Since the influential work of \citet{zoph2017neural},
which formulated NAS as a reinforcement learning problem, the field has
grown rapidly. Modern NAS methods include evolutionary approaches
\citep{real2019regularized}, gradient-based approaches such as DARTS
\citep{liu2019darts}, and one-shot methods \citep{pham2018efficient}.
\citet{ren2024advances} and \citet{ozcelik2026evolutionary} provide
recent surveys.

\modnet{} differs from mainstream NAS in two important respects. First,
the search space is not a discrete set of predefined building blocks
(e.g., convolutional cells or attention modules); it is the space of
all directed graphs over a fixed number of nodes, with node-level
operations drawn from a small set of weighted-sum-plus-threshold units.
Second, the search is performed by a compact evolutionary algorithm
that does not use gradient information, and the resulting architecture
is not a layered stack but a recurrent graph without any modular
structure.

% ------------------------------------------------------------
\section{Summary and Positioning}
\label{sec:related-positioning}
% ------------------------------------------------------------

Table~\ref{tab:positioning} summarizes the relationship between
\modnet{} and the most closely related architectures.

\begin{table}[h]
  \centering
  \small
  \caption{Positioning of \modnet{} relative to related architectures.
    ``Layer-free'' means no predefined layers or grid substrate.
    ``No modules'' means no explicit modular grouping.
    ``Grown'' means both addition and pruning of nodes.
    ``Gradient-free'' means no gradient information at any stage.}
  \label{tab:positioning}
  \begin{tabular}{lccccc}
    \toprule
    \textbf{Architecture} & \textbf{Layer-free} & \textbf{No modules}
      & \textbf{Recurrent} & \textbf{Grown} & \textbf{Gradient-free} \\
    \midrule
    NEAT \citep{stanley2002evolving}           & No  & No  & Optional & Yes & Yes \\
    HyperNEAT \citep{stanley2009hypercube}     & No  & No  & No       & Yes & Yes \\
    Cellular Encoding \citep{gruau1994neural}  & Yes & No  & Yes      & Yes & Yes \\
    Cascade-Correlation \citep{fahlman1990cascade} & No & No & No  & Yes & No  \\
    LNDP \citep{plantec2024evolving}           & Yes & No  & Yes      & Yes & No  \\
    GRN-SO-RC \citep{hajizada2024developmental}& Yes & No  & Yes      & Yes & No  \\
    Neural CA \citep{mordvintsev2020growing}   & Yes & No  & Yes      & Yes & No  \\
    Diff. sorting \citep{petersen2022deep}     & No  & No  & No       & No  & No  \\
    \textbf{\modnet{} (this work)}             & \textbf{Yes} & \textbf{Yes}
      & \textbf{Yes} & \textbf{Yes} & \textbf{Yes} \\
    \bottomrule
  \end{tabular}
\end{table}

The distinguishing features of \modnet{} are:

\begin{enumerate}
  \item \textbf{No predefined layers or substrate geometry.}
  \item \textbf{No modular grouping}---a deliberate departure from our
    own earlier modular variant.
  \item \textbf{Recurrence as a first-class primitive.}
  \item \textbf{Growth through both addition and pruning.}
  \item \textbf{No gradient information used at any stage.}
\end{enumerate}

To the best of our knowledge, no prior architecture simultaneously
satisfies all five of these constraints \emph{and} has been
demonstrated on a benchmark suite of the breadth reported here
(sixteen tasks spanning Boolean logic, digital arithmetic, sorting
networks, and continuous regression).

The remainder of this thesis describes the \modnet{} architecture
(Chapter~\ref{ch:method}), the experimental setup
(Chapter~\ref{ch:experiments}), and the results obtained on the sixteen
benchmark problems (Chapter~\ref{ch:results}).

% ============================================================
%  Bibliography for this chapter
% ============================================================
